\documentclass{handout}

\assignment{Lab 5}
\title{Buoyancy and density}

\begin{document}
\maketitle

\section{Density and Buoyancy}
Go to the Phet simulation called ``Density". Mess around with the `intro' simulation before going on to `compare'.

\begin{questions}
\question What determines the amount of water displaced when an object is thrown into the water? Try to find both a block that floats and a block that sinks.

\answerspace{1.5in}

\question Go to the compare version: if all blocks have the same mass, what determines if they sink or float?

\answerspace{1.5in}

\question If all blocks have the same volume, what determines if they sink or float?

\answerspace{1.5in}

\question Go to the `mystery' section of the simulation, and find the object in set 2 made of glass. Explain your reasoning.
\end{questions}

\section{Measuring the density of two objects}
We will use a force sensor with PASCO capstone to measure the density of two objects of unknown material.
\begin{questions}
    \question What are the masses of your two objects?

    \answerspace{1.5in}

    \question What are the volumes of your two objects? Explain your answers

    \answerspace{\fill}
    \clearpage

    \question When submerged in water, what are the `apparent' masses of your two objects?

    \answerspace{1.5in}

    \question Using this information, find the density of the two objects. Explain  your reasoning.

    \answerspace{\fill}

    \question What materials do you think your objects are made of?

    \answerspace{2in}

\end{questions}

\end{document}
